\documentclass[10pt,a4paper]{amsart}
\usepackage[margin=19mm]{geometry}
\usepackage{amsmath,amssymb,mathtools}
\usepackage[scr=rsfs,cal=euler]{mathalfa}
\usepackage{xfrac}
\usepackage[unicode,hidelinks,bookmarks=false]{hyperref}
\newcommand{\ie}{i.e.\ }
\newcommand{\cf}{cf.\ }
\newcommand{\resp}{respectively\ }
\newcommand{\Subsection}{\subsection}
\providecommand{\nobreakdash}{\nobreak\hbox{-}}
\setlength{\emergencystretch}{2em}
\multlinegap0pt
\usepackage{bm}
\usepackage{bbm}
\usepackage{dsfont}
\usepackage{xspace}
\usepackage{cancel}
\usepackage{mathtools}
\usepackage{amsthm}
\makeatletter
\@ifundefined{prop}{\newtheorem{prop}{Proposition}}{}
\makeatother
\title[The sharp $C^0$-fragmentation property for Hamiltonian diffe]{The sharp $C^0$-fragmentation property for Hamiltonian diffeomorphisms and homeomorphisms on surfaces}
\author{Baptiste Serraille}
\date{Source: arXiv:2403.15767v2 (CC BY 4.0)}
\begin{document}
\maketitle

\noindent\emph{Source and licence.} The sharp $C^0$-fragmentation property for Hamiltonian diffeomorphisms and homeomorphisms on surfaces. Version arXiv:2403.15767v2 (CC BY 4.0), distributed under the Creative Commons Attribution 4.0 International licence, as recorded in the arXiv licence metadata for this exact version and shown on the abstract page https://arxiv.org/abs/2403.15767v2. Full rights notice: licenses/arxiv-2403.15767-CC-BY-4.0.txt.\par\smallskip

\noindent\textit{Editorial scope.} Counted unit: the Prop whose source label is \texttt{prop: adjusting volume forms} in the article (source0.tex lines 673--682, source label \texttt{prop: adjusting volume forms}), delivered together with the article's own complete original proof of it (source0.tex lines 684--717, one \texttt{proof} environment of 34 lines). No mathematics is written, repaired or re-derived: statement and proof are the article's own text line for line. Required context retained uncounted: prop: Moser's trick (lines 650--669). Every definition, display and label of those lines is retained unchanged. Omitted from this package and not redistributed: 1--649, 670--672, 683--683, 718--967. Nothing inside a retained excerpt is omitted or altered: every retained body is delivered exactly as the article's lines are written, so no omission receipt of any kind accompanies this package. Citations: the retained excerpts carry no citation command at all, so no bibliography entry of the article is transcribed and no reference is redistributed. Reference adaptation: none was needed and none was made; every reference inside the delivered unit resolves inside it, so no unresolved reference or citation is produced. Printed slips kept literal: no printed word, symbol, hypothesis or number of the counted statement or of its proof was altered. Layout and class: the article's own class is \texttt{amsart} with packages including geometry, xcolor, hyperref, tikz, graphicx, amssymb, eucal, mathrsfs, latexsym, amsmath, epsfig, epic, eepic, amscd; the wrapper is the lane's pinned \texttt{amsart} editorial wrapper at 10pt on A4, which restates the theorem-like environments the retained text needs and adds the article's own macro definitions that the retained text needs (none), with extra wrapper packages (bm, bbm, dsfont, xspace, cancel, mathtools). The delivered numbering is the unit's own. Assets: no retained span contains a figure environment, a graphics-inclusion command, a PDF-page inclusion, a table or a TikZ picture; no image file is copied into this package, the package declares no asset dependency and no image file is redistributed with it. Licence: version arXiv:2403.15767v2 of this article is distributed under the Creative Commons Attribution 4.0 International licence, as recorded in the arXiv licence metadata for this exact version and shown on the abstract page https://arxiv.org/abs/2403.15767v2. A full rights notice is delivered at licenses/arxiv-2403.15767-CC-BY-4.0.txt. Characters: every retained excerpt is pure ASCII, so the delivered engine, XeLaTeX with its default Latin Modern text font, reports no missing character.
\begin{prop}[Moser's trick]\label{prop: Moser's trick}
Let~$M$ be a compact connected oriented manifold of dimension~$n$, possibly with a non-empty boundary~$\partial M$, and let~$\omega_1$,~$\omega_2$ be two volume forms on~$M$. Assume that~$\int_M \omega_1=\int_M\omega_2$. If~$\partial M\neq 0$, we also assume that the forms~$\omega_1$ and~$\omega_2$ coincide on~$\partial M$.

Then there exists a diffeomorphism~$f: M \rightarrow M$, isotopic to the identity, such that~$f^* \omega_2=\omega_1$. Moreover,~$f$ can be chosen to satisfy the following properties: 

\medskip


(i) If~$\partial M \neq \emptyset$, then~$f$ is the identity on~$\partial M$, and if~$\omega_1$ and~$\omega_2$ coincide near~$\partial M$, then~$f$ is the identity near~$\partial M$.

\medskip
 
(ii) If~$M$ is partitioned into polyhedra (with piece-wise smooth boundaries), so that~$\omega_1-\omega_2$ is zero on the~$(n-1)$-skeleton~$\Gamma$ of the partition and the integral of~$\omega_1$ and~$\omega_2$ on each of the polyhedra are equal, then~$f$ can be chosen to be the identity on~$\Gamma$.
 
\medskip
 
(iii) Suppose that~$\omega_2=\chi\omega_1$ for a function~$\chi$, then the diffeomorphism~$f$ can be chosen to satisfy the following estimate: 
$$\Vert f \Vert_{C^0}\leq C \lVert\chi-1 \rVert_{C^0},$$
for some~$C>0$. Here,~$\lVert \cdot \rVert_{C^0}$ denotes the standard sup norm on functions.
\end{prop}

\begin{prop}\label{prop: adjusting volume forms}
Let~$C\in \mathbb{R}$ be a constant and~$M^{n-1}$ a compact manifold equipped with a volume form~$\omega'$ and a distance~$d'$, we also define~$d$ the product distance on~$M^{n-1}\times [-1,1]$. Let~$\omega=\omega' \wedge dz$ and~$\Omega$ two volume forms on~$M^{n-1}\times [-1,1]$, where~$z$ denote the coordinate on~$[-1,1]$. Let~$\chi$ be the function such that~$\omega=(1+\chi)\Omega$. We assume that: 

\begin{itemize}
    \item~$\int_{M^{n-1}\times [-1,1]}\omega =\int_{M^{n-1}\times [-1,1]}\Omega.$
    \item~$\Vert \chi \Vert_{C^0}\leq C.$
    \item There exists~$\delta>0$ such that~$\text{Supp}(\chi)\subset M^{n-1}\times [0,\delta].$
\end{itemize}
Then, there exists a constant~$D\in \mathbb{R}$ independent of~$\delta$ such that we can find~$f\in \text{Diff}(M^{n-1}\times[-1,1])$,~$f\equiv Id$ on~$M^{n-1}\times [-1,0]$,~$f^*\Omega=\omega$ and~$\Vert f \Vert_{C^0}\leq D\delta$.
\end{prop}

\begin{proof}
We will denote by~$z$ the last coordinate of the manifold~$M^{n-1}\times [-1,1]$ and by~$x$ the coordinate on~$M^{n-1}$. Let us describe two diffeomorphisms of~$N^n:=M^{n-1}\times \mathbb{R}$, which will be used in the definition of the desired map $f$ on~$M^{n-1}\times [-1,1]$.

Define~$\rho_\delta: \mathbb{R}\rightarrow \mathbb{R}, x\mapsto 2\delta x$. We can then define two maps from~$N^n$ to~$N^n$ by:
$$\Psi_1(x,z):=\left(x,\int_0^z(1+\chi(x,t))\,dt\right)$$
and 
$$\Psi_2(x,z):=\left(x,\int_0^z(1+\rho_\delta'(t)\chi(x,\rho_\delta(t)))\,dt\right).$$

We can then compute 
$$d\Psi_1(x,z):=\begin{pmatrix}
	Id & *\\
	0 &  1+\chi(x,z)
\end{pmatrix},$$
and
$$d\Psi_2(x,z):=\begin{pmatrix}
	Id & *\\
	0 &  1+\rho_\delta'\chi(x,\rho_\delta(z))
\end{pmatrix}.$$

Since~$1+\chi>0$ the function~$\Psi_1$ is a diffeomorphism, in the same manner~$\Psi_2$ is also a diffeomorphism. A simple computation also shows that~$(\Psi_{1})^*\omega=\text{det}(d \Psi_1)\omega=(1+\chi(x,z))\omega=\Omega$ and also~$(\Psi_{2})^*\omega=(1+\rho_\delta'\chi(x,\rho_\delta(z)))\omega$ which gives 
$$\left(((\Psi_1)^{-1} \circ \Psi_2)^*\Omega\right)_{(x,z)}=(1+\rho_\delta'(z)\chi(x,\rho_\delta(z)))\omega_{(x,z)}.$$

Moreover, 
$$d\left((x,z),\Psi_1(x,z)\right)=\left\vert\int_0^z \chi(x,\rho_\delta(t))\,dt\right\vert \leq \Vert \chi \Vert_{C^0} \delta\leq C \delta$$
so~$\Vert \Psi_1 \Vert_{C^0}\leq C\delta$ and similarly~$\Vert \Psi_2 \Vert_{C^0}\leq C \delta$.

If~$z>\delta$, since~$\chi$ has support in a strip we have~$\Psi_1(x,z)=(x,z+c(x))$, where~$c(x)$ is a function independent of~$z$ and whose value is~$c(x)=\int_0^\delta \chi(x,t)\, dt$. If~$z>0.5$, we have, by the same argument, that~$\Psi_2(x,z)=(x,z+d(x))$, where~$d(x)=\int_0^{0.5}\rho_\delta'(t)\chi(x,\rho_\delta(t))\, dt$. 

It follows that those two diffeomorphisms aren't compactly supported in~$M^{n-1}\times [-1,1]$, however since~$c(x)=d(x)$ (simple change of variable in the integral),~$\Psi:=(\Psi_1)^{-1} \circ \Psi_2$ can be restricted on~$M^{n-1}\times [-1,1]$ to a compactly supported diffeomorphism. Moreover,~$$\left(\Psi^*\Omega\right)_{(x,z)}=\left(((\Psi_1)^{-1} \circ \Psi_2)^*\Omega\right)_{(x,z)}=(1+\rho_\delta'(z)\chi(x,\rho_\delta(z)))\omega_{(x,z)}$$ and 
$$\Vert\Psi\Vert_{C^0}\leq \Vert \Psi_1\Vert_{C^0}+\Vert \Psi_2\Vert_{C^0}\leq 2C\delta.$$

By Proposition \ref{prop: Moser's trick} one can find a function~$h\in \text{Diff}_c(M^{n-1}\times[-1,1])$ such that~$h^*(\Psi^*(\Omega))=\omega$ and there exists a constant~$C'$ independent of~$\omega$ and~$\Omega$ such that~$\Vert h \Vert_{C^0}\leq C' \Vert \rho_\delta' \chi(\cdot,\rho_\delta(\cdot))\Vert_{C^0}\leq 2C'C\delta$. It then follows that~$\Psi h$ is the diffeomorphism we needed.

\end{proof}

\end{document}
